In this work, we introduce \textbf{Doob Policy Optimization (DoobPO)}, a unified framework for RL fine-tuning of diffusion and flow policies.
The framework supports a broad family of policy updates and avoids exact action-likelihood evaluation by separating policy improvement from generative policy learning.
Theoretically, we justify DoobPO from two directions.
First, we characterize the ratio-tilted diffusion process through Doob's $h$-transform
(\cref{lem:algo_correctness}).
We then show that, for each fixed state, the gradient of the conditional drift matching objective is proportional to that of the ideal objective for both diffusion and flow policies
(\cref{thm:grad_equi_diff,thm:grad_equi_fm}).
Second, under the stated assumptions, the population minimizer of conditional drift matching induces a diffusion process that solves the corresponding endpoint-constrained stochastic optimal control problem
(\cref{prop:cdm_solves_soc}).
The resulting process is the minimum-energy modification of the reference diffusion that realizes the prescribed tilted endpoint marginals.
This provides a control-theoretic interpretation of how the desired policy update is transferred into generative dynamics.
Experimentally, we evaluate DoobPO on MuJoCo benchmarks with both diffusion and flow policies and observe superior overall performance against strong baselines
(\cref{sec:diffusion_exp,sec:flow_exp}).

More broadly, DoobPO provides a common interface between policy optimization and generative model training.
Policy updates with analytic ratios and updates represented by learned ratio models are accommodated within the same conditional drift matching framework
(\cref{sec:unified_framework}).
This separates the choice of the policy update rule from the choice of the generative policy class.
Nevertheless, our study mainly focuses on continuous generation, and the empirical evaluation is limited to MuJoCo continuous-control benchmarks and Robomimic manipulation tasks.
The theoretical results characterize population objectives under the stated assumptions.
An important direction for future work is to quantify how ratio estimation error and approximate drift matching affect the realized policy update.
We also plan to extend the framework to discrete generation and evaluate its applicability to more complex policy learning settings.